\scanpage{74}
Bien entendu, $Q(A)$ est un $\handwritten{\lambda}$-anneau, et sa $\handwritten{\lambda}$-structure se lit sur la structure d'anneau, puisqu'on aura
\[
\handwritten{\lambda}_t(\handwritten{\varphi}_a)=1+\handwritten{\varphi}_a t\ ,
\]
D'autre part, un homomorphisme d'anneaux augmentés $Q(A)\handwritten{\to}\underline{Z}\times(1+\overset{\handwritten{\scriptscriptstyle\wedge}}{A}^{\handwritten{+}})$ est automatiquemejt un $\handwritten{\lambda}$-homomorphisme. Je ne crois pas qu'il y ait lieu de dire plus des $\handwritten{\lambda}$-structures (sauf peut-être en remarque que la filtration mise sur $Q(A)$ est celle déduite de sa $\handwritten{\lambda}$-structure).

2) \uline{Cas où} $A$ \uline{est un anneau local}, de corps résiduel $k$ tel que $\mathrm{car}\ k\handwritten{\ne}2$. Il semble bien que tout ce qui précède marqhe dans ce cas; bien entendu, $H^{\handwritten{\bullet}}(A)$ désigne la cohomologie du ``groupe fondamental'' de $\mathrm{Spec}(A)$, au sens de mon séminaire. Il faudrait alors tout rédiger dans le cadre: formes quadratiques sur un anneau local. N.B. une telle forme est dite non dégénérée si la réduite est non dégénérée; en principe, il \deleted{\unsure{faudrait être ca|faudrait être ça}} devrait être possible de façon assez triviale de relever les constructions habituelles sur $k$, en des constructions sur $A$.

Bien entendu, il serait très agréable de dépasser le programme minimum et d'arriver à une théorie des classes de Stiefel-Whitney pour des Modules quadratiques sur un \inserted{pré}schéma quelconque $S$, \deleted{\unsure{On pourrait commencer|On pourrait commen}} en supposant au besoin que les caractéristiques résiduelles sont $\handwritten{\ne}2$. \deleted{trit} A titre de test, on peut toujours essayer de prouver l'invariance de la définition triviale relative à un Module\inserted{ quadratique} qui splitte en modules\inserted{ quadratiques} de rang 1.

Comme prime au lecteur, on peut en tous cas donner la compatibilité de la définition purement algébrique des classes de Stiefel-Whitney, et de la définition topologique, lorsque $A$ est l'anneau local\inserted{ d'un point $s$} d'un schéma de type fini\inserted{$S$} sur $\underline{C}$. Un module quadratique sur $A$ définit un module quadratique sur un voisinage de $S$, donc $S$ si on veut, i.e. un fibré prinnipal homogène algébrique sous le schéma en groupes orthogonal. Par Gaga il lui correspond un fibré principal analytique, donc topologique, d'où des classes de Stiefel-Whitney topologiques dans $H^{\handwritten{\bullet}}(S_{\mathrm{top}},\underline{Z}/2\underline{Z})$. Or on a un homomorphisme canonique,que tu définiras:
\[
H^{\handwritten{\bullet}}(S_{\mathrm{top}},\underline{Z}/2\underline{Z})\handwritten{\longrightarrow}H^{\handwritten{\bullet}}(A,\underline{Z}/2\underline{Z})\ ,
\]
et on voit aussitôt que les images des classes caractéristiques topologiques sont les classes caractéristiques algébriques, (car il suffit de le vérifier en rang 1).
% archive sheet number: 73
